\documentclass[12pt,a4paper]{article}
\usepackage[utf8]{inputenc}
\usepackage[T1]{fontenc}
\usepackage[brazil]{babel}
\usepackage{amsmath,amssymb,amsthm}
\usepackage{geometry}
\geometry{margin=2.5cm}

\newtheorem{theorem}{Teorema}[section]
\newtheorem*{proofsketch}{Demonstração}

\title{Teorema de Pitágoras}
\author{}
\date{}

\begin{document}
\maketitle

\begin{theorem}[Pitágoras]
Em todo triângulo retângulo, o quadrado da medida da hipotenusa é igual
à soma dos quadrados das medidas dos catetos. Em símbolos: se $a,b$ são
os catetos e $c$ é a hipotenusa, então
\[
  a^{2} + b^{2} = c^{2}.
\]
\end{theorem}

\begin{proof}
Seja $ABC$ um triângulo retângulo com ângulo reto em $C$, catetos
$BC=a$, $AC=b$ e hipotenusa $AB=c$. Tracemos a altura $CH$ relativa
à hipotenusa, com $H \in AB$, e denotemos $AH=m$ e $BH=n$, de modo que
$m+n=c$.

Os triângulos $ACH$ e $ABC$ são semelhantes, pois compartilham o ângulo
em $A$ e ambos possuem um ângulo reto (em $H$ e em $C$, respectivamente).
Da semelhança, a razão entre o cateto adjacente a $A$ e a hipotenusa
é a mesma nos dois triângulos:
\[
  \frac{AH}{AC} = \frac{AC}{AB}
  \quad\Longrightarrow\quad
  \frac{m}{b} = \frac{b}{c}
  \quad\Longrightarrow\quad
  b^{2} = mc. \tag{1}
\]

Analogamente, os triângulos $CBH$ e $ABC$ são semelhantes (ângulo comum
em $B$ e ambos retângulos). Logo,
\[
  \frac{BH}{BC} = \frac{BC}{AB}
  \quad\Longrightarrow\quad
  \frac{n}{a} = \frac{a}{c}
  \quad\Longrightarrow\quad
  a^{2} = nc. \tag{2}
\]

Somando (1) e (2):
\[
  a^{2}+b^{2} = nc + mc = c(m+n) = c\cdot c = c^{2}.
\]
Isso conclui a demonstração.
\end{proof}

\end{document}
